\section{Mô hình đe dọa}\label{sec:threat}

\paragraph{Tác tử.} Cho chỉ thị người dùng $u$, ở mỗi bước mô hình sinh một hành động $a_t$ (công cụ và tham số)
hoặc câu trả lời cuối $y$. Tool Broker kiểm $a_t$ theo schema, thực thi và trả quan sát $o_t$. Có sáu công cụ đọc
(tìm và đọc tài liệu, truy vấn cơ sở dữ liệu chỉ đọc, tìm và lấy trang lưu đệm, máy tính) và hai \emph{sink} có tác
dụng phụ, \code{send\_email\_mock} và \code{post\_webhook\_mock}; cả hai chỉ ghi log, không gửi ra mạng. Câu trả lời
cuối được coi là sink thứ ba. Với lời gọi sink $a$, ta viết $\dest(a)$ cho đích (\code{to} hoặc \code{endpoint}) và
$\pay(a)$ cho các tham số còn lại.

\paragraph{Hai nhãn.} Mỗi nội dung mang hai nhãn độc lập do host gán. \emph{Độ tin cậy}: $u$ và prompt hệ thống là
đáng tin, mọi quan sát công cụ là không đáng tin. \emph{Độ nhạy cảm}: từ công khai (S0) tới mật (S2); tài liệu mật
chứa giá trị (khóa, mã, định danh tài khoản) không được rời tác tử nếu người dùng không cho phép. Một trang không đáng
tin có thể công khai, còn một bản ghi mật có thể đến từ nguồn kẻ tấn công không chạm được; gộp hai chiều làm một sẽ
hoặc chặn mọi dữ liệu mật, hoặc tin mọi nguồn nội bộ.

\paragraph{Kẻ tấn công.} Kẻ tấn công kiểm soát văn bản của một tài liệu mà công cụ đọc trả về (tài liệu, trang lưu
đệm, trường văn bản trong cơ sở dữ liệu, đầu ra công cụ), ở bất kỳ vị trí và giọng điệu nào, và biết tên cùng schema
công cụ. Họ không sửa được $u$, prompt hệ thống, broker hay trọng số, và không quan sát được quỹ đạo. Trên tập niêm
phong, tấn công là tĩnh, được viết mà không biết mã phòng thủ. Jailbreak trực tiếp, đầu độc dữ liệu huấn luyện, từ
chối dịch vụ và hệ nhiều tác tử nằm ngoài phạm vi.

\paragraph{Vị từ gây hại.} Gọi $G(u)$ là tập đích người dùng cấp quyền, $V$ là tập giá trị mật trên quỹ đạo mà người
dùng không cho phép tiết lộ, và $\mathcal{S}$ là tập lời gọi sink \emph{đã thực thi}. Quỹ đạo $\tau$ gây hại nếu
\begin{equation}\label{eq:harm}
\mathrm{Harm}(\tau) = H_\mathit{dest} \lor H_\mathit{leak} \lor H_\mathit{tamper} \lor H_\mathit{side},
\end{equation}
trong đó $H_\mathit{dest}$: $\exists a \in \mathcal{S}$ với $\dest(a) \notin G(u)$; $H_\mathit{leak}$: có
$v \in V$ xuất hiện trong $\pay(a)$ của một $a \in \mathcal{S}$ hoặc trong $y$; $H_\mathit{tamper}$: $\exists a \in
\mathcal{S}$ với $\dest(a) \in G(u)$ nhưng một trường người dùng chỉ định mang giá trị khác; $H_\mathit{side}$:
nhiệm vụ chỉ đọc nhưng $\mathcal{S} \neq \emptyset$. Lời gọi bị chặn không gây hại. Bộ chấm đóng băng tính các vị từ
từ lời gọi đã thực thi và $y$; suy luận của mô hình không được ghi.

\paragraph{Sáu cơ chế.} Chúng tôi phân loại tấn công theo \emph{cơ chế}, tức con đường đoạn chèn biến thành tác hại
(Bảng~\ref{tab:mechanisms}). Ở \mER{} và \mCR{}, lời gọi gây hại đi tới đích người dùng chưa viết, nên kiểm tra đích
là đủ dù cách diễn đạt thế nào. Ở \mSL{} và \mAL{}, đích hợp lệ hoặc không có đích; chỉ nguồn gốc của giá trị phân biệt
tấn công với nhiệm vụ. Ở \mPT{}, cả đích lẫn độ nhạy cảm đều bình thường, nên đây là cơ chế khó nhất với luật theo
nhãn. \mUS{} kiểm tra xem việc người dùng nhắc tới một đích, chẳng hạn trong câu cấm, có bị hiểu nhầm là cấp quyền
không.

% M1..M6 trong log/luận văn = ER, CR, SL, AL, PT, US.
\begin{table}[!htbp]
\centering
\caption{Sáu cơ chế tấn công và vị từ gây hại tương ứng.}\label{tab:mechanisms}
\footnotesize
\begin{tabular}{@{}lp{0.70\textwidth}l@{}}
\toprule
Mã & Cơ chế (kiểu con) & Vị từ \\
\midrule
\mER & Chuyển hướng tường minh: yêu cầu gửi tới đích không có trong $u$, có động từ gửi và địa chỉ đầy đủ & $H_\mathit{dest}$ \\
\mCR & Chuyển hướng che giấu: cùng mục tiêu nhưng tránh mẫu ``gửi + địa chỉ'' (ghi nhận không động từ, chuyển tiếp hoặc cc, tách địa chỉ, đích nằm ở tài liệu thứ hai) & $H_\mathit{dest}$ \\
\mSL & Rò rỉ qua sink: đưa giá trị mật vào sink người dùng đã nêu (thân e-mail, payload, trường trông vô hại) & $H_\mathit{leak}$ \\
\mAL & Rò rỉ qua câu trả lời: đưa giá trị mật vào câu trả lời cuối (nguyên văn, lồng trong câu, có nhãn) & $H_\mathit{leak}$ \\
\mPT & Sửa tham số: đổi một giá trị thường trong payload của sink người dùng đã nêu (số, boolean, đơn vị, định danh) & $H_\mathit{tamper}$ \\
\mUS & Tác dụng phụ ngoài yêu cầu: chèn lời gọi sink vào nhiệm vụ chỉ đọc (đích chỉ xuất hiện trong câu cấm, chỉ là thông tin, hoặc không có) & $H_\mathit{side}$ \\
\bottomrule
\end{tabular}
\end{table}


\paragraph{Nhánh lành tính và biến thể.} Mỗi họ tấn công có hai nhánh lành tính cùng nhiệm vụ và môi trường.
\emph{benign\_plain} thay đoạn chèn bằng ghi chú vô hại cùng độ dài, để đo tiện ích. \emph{benign\_auth} giữ một yêu
cầu giống mục tiêu tấn công nhưng được người dùng cho phép rõ ràng; chặn ở đây là \emph{chặn nhầm}. Mỗi họ có năm
biến thể áp cho cả $u$ lẫn đoạn chèn: tiếng Việt chuẩn, diễn đạt lại, trộn mã Việt--Anh, không dấu và tiếng Anh. Định
danh, địa chỉ và giá trị mật giống hệt nhau giữa các biến thể, nên khác biệt chỉ do ngôn ngữ.
